\section{Introducción al IoT y Sensores}
En el ámbito del Internet de las Cosas (IoT, Internet of Things), los sensores juegan un papel fundamental al permitir captar variables del entorno físico y comunicarlas a sistemas conectados para su procesamiento y análisis \cite{webology2018}. Esto convierte a los sensores en elementos clave para la automatización, el monitoreo y la toma de decisiones en tiempo real.
\\\\
Para el monitoreo de temperatura, se han utilizado ampliamente sensores económicos como el \textbf{LM35} y el \textbf{DHT11}, gracias a su bajo costo y facilidad de integración con microcontroladores. Por ejemplo, Prasad y Jahnavi (2018) desarrollaron un sistema IoT en el que un sensor LM35 conectado a un Arduino permite monitorear la temperatura de una habitación y enviar los datos a la nube para su visualización remota \cite{webology2018}. De manera similar, se han reportado implementaciones donde el sensor DHT11, acoplado a módulos inalámbricos como el ESP8266 o el ESP32, mide simultáneamente temperatura y humedad, y transmite los datos a una plataforma en línea para su análisis y almacenamiento \cite{webology2018}.
\section{Comunicación de Datos: El Rol del Protocolo MQTT}
En cuanto al envío de datos de sensores a través de internet, el protocolo de mensajería \textbf{MQTT} (Message Queuing Telemetry Transport) se ha consolidado como uno de los estándares predominantes en IoT, tal como lo describe la documentación de Amazon Web Services \cite{awsMQTT}. Este protocolo fue diseñado para ser ligero y eficiente, lo que lo hace ideal para redes con recursos limitados y dispositivos embebidos que operan con conectividad inalámbrica.
\\\\
Actualmente, MQTT es ampliamente adoptado en la industria para aplicaciones como telemetría remota, automatización del hogar, monitoreo agrícola, ciudades inteligentes, entre otras. Su arquitectura basada en el modelo \textit{publish/subscribe} permite que los dispositivos (nodos sensores) publiquen datos en un \textit{broker} central, donde otros clientes suscriptores pueden acceder a esa información en tiempo real. Esta estructura mejora la escalabilidad, la modularidad y la eficiencia en comparación con los esquemas cliente-servidor tradicionales \cite{awsMQTT}.
\\\\
Dispositivos como el ESP32 o el Arduino Mega son frecuentemente utilizados para este propósito, ya que permiten leer datos desde sensores y publicarlos periódicamente al broker MQTT utilizando módulos Wi-Fi o Ethernet.
\section{Expansión de Variables Sensadas y Hardware Compatible}
Además de la temperatura y la humedad, los sistemas IoT pueden incorporar sensores para medir una variedad de variables como presión barométrica, concentración de gases ($CO_2$, metano, monóxido de carbono), niveles de luz, vibración, o incluso parámetros biométricos en aplicaciones de salud. Según Prasad y Jahnavi (2018), la incorporación de sensores adicionales permite ampliar la funcionalidad de los sistemas y adaptarlos a contextos específicos como agricultura de precisión, sistemas de alerta temprana o control industrial \cite{Prasad2018}.
\\\\
Placas de desarrollo como la Raspberry Pi o el STM32 ofrecen una mayor capacidad de procesamiento local, lo cual es útil para aplicaciones que requieren preprocesamiento de datos, agregación o ejecución de modelos de inteligencia artificial antes de enviar la información a la nube.
\section{Plataformas de Visualización y Procesamiento}
El almacenamiento, procesamiento y visualización de los datos recolectados es un componente esencial en cualquier sistema IoT. Plataformas como \textbf{ThingSpeak}, \textbf{Node-RED} y \textbf{Blynk} permiten diseñar paneles de control personalizados, generar alertas automáticas y realizar análisis en tiempo real de las variables monitoreadas. Según la documentación de AWS (2023), estas herramientas proporcionan interfaces gráficas intuitivas, integraciones con protocolos como MQTT y opciones para la ejecución de scripts o flujos lógicos \cite{AWS-MQTT}.
\\\\
Estas plataformas también permiten implementar reglas de decisión, exportación de datos históricos y conexión con servicios de terceros como Google Sheets, IFTTT o servicios de mensajería, lo cual incrementa significativamente la funcionalidad del sistema sin requerir una infraestructura compleja.
\section{Consideraciones de Seguridad}
Uno de los desafíos más relevantes en los sistemas IoT actuales es la seguridad. La transmisión constante de datos a través de internet, combinada con dispositivos de bajo costo y capacidades limitadas, representa un punto crítico para la integridad y confidencialidad de la información. Por ello, se recomienda implementar cifrado en las comunicaciones MQTT (por ejemplo, usando TLS/SSL), autenticación de dispositivos, y mecanismos de detección de intrusiones. Aunque no todos los sensores o microcontroladores soportan estas medidas de forma nativa, su implementación se vuelve indispensable en aplicaciones sensibles como monitoreo industrial o médico.